\section{Residuals in the dual energy norm}\label{sec:energy}

We now bound the two residuals from Lemma~\ref{lem:duhamel} against an
arbitrary test vector, using the one-period energy of that vector.
The scalar estimate of Section~\ref{sec:scalar} is not sufficient for
this purpose: the test vector need not resemble the free potential.
Instead, we represent propagated sources as divergences of flows and
pair those flows with the gradient of the propagated test vector.
For a group of sources associated with one obstacle strand, we first
route the signed mass to a nearby slot.  Only the resulting net mass
must then be sent to infinity.  The cancellation in
Lemma~\ref{lem:cancellation} reduces this second cost.

\begin{proposition}[Dual energy estimate]\label{prop:dualenergy}
Fix \(S\), and use the scales in \eqref{eq:scales}.  Suppose that
\(B_{50R_*}(x)\) lies in the interior of the finite cube and that no
point of \(E\) lies above this ball.  Let \(\mathcal G\) be the good
event of Proposition~\ref{prop:good}, and let \(\mu_+,\mu_-\) be the
terminal residuals in \eqref{eq:residualdefinition}.
There are nonnegative measurable functions \(X_+,X_-\) of the exposure
and a deterministic rapidly small function \(\rho(\beta)\) such that,
on \(\mathcal G\),
\begin{equation}\label{eq:dualenergy}
 \abs{\ip{\mu_d}{y}}
 \le X_d\sqrt{e_y}+C_0\rho(\beta),
 \qquad d\in\{+,-\},
\end{equation}
for every real cut vector \(y\) that vanishes on the cut killing slots
and satisfies \(\norm{y}_\infty\le C_0\).  The functions \(X_d\) do not
depend on \(y\) or \(C_0\), and
\begin{equation}\label{eq:dualenergyexpectation}
 \EE_P\!\left[\ind_{\mathcal G}(X_+^2+X_-^2)\right]
 =o(\beta^{-5/2}).
\end{equation}
All estimates are uniform over the finite cubes, deterministic pin sets,
and injection slots satisfying the stated assumptions.
\end{proposition}

We work in one direction at a time, writing \(U_d=U\) for the forward
direction and \(U_d=U^*\) for the reverse direction.  Time in the
following argument is elapsed time in this chosen traversal; all
instantaneous operations retain the ordering fixed in
Sections~\ref{sec:loops} and~\ref{sec:sources}.
Write \(T_{t,s}\) for its propagator from \(s\) to \(t\), and extend
this notation periodically when \(t>\beta\).

For a terminal vector \(y\), the adjoint suffix
\(Y(t)=T_{\beta,t}^*y\) satisfies, by Lemma~\ref{lem:kernelenergy},
\begin{equation}\label{eq:energy-suffix}
 \int_0^\beta\D_t(Y(t))\,dt\le e_y,
 \qquad
 \norm{y-U_d^*y}_2^2\le2e_y,
 \qquad
 \norm{Y(t)}_\infty\le C_0.
\end{equation}
These bounds hold for signed vectors.  In particular, they require
neither a reversible one-period kernel nor time-reversal invariance of
the exposure law.

At a fixed slice, a finite-energy flow \(J\) with divergence \(\nu\)
gives
\begin{equation}\label{eq:flowpairing-energy}
 \abs{\ip{\nu}{Y}}
 =\left|\sum_eJ(e)\bigl(Y(e^-)-Y(e^+)\bigr)\right|
 \le\bigl(2\norm{J}_2^2\D_t(Y)\bigr)^{1/2}.
\end{equation}
Here \(Y\) is extended by zero on the killed inner boundary and outside
the cube.  It has finite support, so summation by parts has no term at
infinity, and boundary-crossing edges are already included in
\(\D_t\).  Nonzero total mass is allowed: the flows in
Lemma~\ref{lem:flowoverlap} carry that mass to infinity on the extended
available-slot graph.  This is the same flow-pairing principle as in
\cite[Section~3.3]{BlochLaw}; the block cancellation below provides the
additional saving needed here.

\subsection{Flows for the sources of one block}

Divide the sources in the chosen direction into two collections.
The first contains all cutoff sources, all obstacle sources from the
smooth potential, and the point-profile obstacle sources whose elapsed
times lie in \([\beta/2,\beta]\).  The second contains the remaining
point-profile obstacle sources, at elapsed times in \([0,\beta/2)\).
We begin with the first collection.  Its obstacle sources satisfy the
block bounds of Lemma~\ref{lem:cancellation} with envelope \(W\),
and its cutoff sources have size at most \(Cb_*\) per slot and unit
time, where
\[
 b_*=\beta^{-1}R^{-3}
\]
is the bound in \eqref{eq:backgroundsource}.  Since \(\beta/2\) is a
block division, no block is truncated when the late point-profile
sources are included.

Let \(I_1,\ldots,I_n\) be the consecutive blocks in this direction,
with lengths between \(L\) and \(2L\).  Thus \(n\le C\beta/L\).
For each block \(I_j\), let \(\mathcal A_j\) be the obstacle slots at
its starting slice whose sites lie in \(B_{4R}(x)\), and use the same
index for the strand starting at such a slot.  Denote its starting
site by \(z_\alpha\).  Set
\[
 \theta_\alpha=
 \begin{cases}
 L^{-1/4},&\text{if this strand is isolated in }I_j,\\
 1,&\text{otherwise},
 \end{cases}
 \qquad
 Q_j=\sum_{\alpha\in\mathcal A_j}W(z_\alpha)^2,
\]
and
\[
 C_j=\sum_{\alpha,\alpha'\in\mathcal A_j}
 \frac{\theta_\alpha\theta_{\alpha'}
 W(z_\alpha)W(z_{\alpha'})}
 {1+|z_\alpha-z_{\alpha'}|_\infty}.
\]
These are finite snapshot sums, with slot multiplicities retained.

\begin{lemma}[One-block flow bound]\label{lem:blockflow}
For the first collection of sources, fix a block \(I_j\) on
\(\mathcal G\).  Propagate its sources to a fully regular slice \(t\)
of \(I_{j+1}\), interpreting \(I_{n+1}=\beta+I_1\) in a copy of the
period.  Their propagated signed mass differs in counting
\(\ell^1\) norm by a rapidly small amount from a localized mass
\(\nu_j(t)\), supported in \(B_{6R}(x)\), which is the divergence
of a flow of squared norm at most
\begin{equation}\label{eq:blockflowbound}
 F_j=C\left(L^6Q_j+L^2C_j+L^2b_*^2R^5\right).
\end{equation}
The constant and the localization error are uniform in \(j\) and
\(t\).  Moreover,
\begin{align}\label{eq:blockflowexpectation}
 \EE_P[\ind_{\mathcal G}F_j]
 \le C\bigl[&L^6\lambda\beta^{-2}D_*^{-1}
 +L^2(L^{-1/2}\lambda^2+L^3\lambda^3)\beta^{-2}R
 +L^2b_*^2R^5\bigr].
\end{align}
If the sources of the final block are propagated only to the terminal
cut, their mass differs by a rapidly small amount in \(\ell^1\)
from a mass \(m_{\mathrm{loc}}\) satisfying
\begin{equation}\label{eq:lastblocknorm}
 \norm{m_{\mathrm{loc}}}_2^2
 \le K_n:=CL^5(Q_n+b_*^2R^3).
\end{equation}
\end{lemma}

\begin{proof}
Every obstacle source in the block is associated with a strand starting
within \(B_{4R}(x)\).  Sources of that strand are within \(CL\) of
its start, and have total variation at most \(CLW(z_\alpha)\).
Lemma~\ref{lem:localization} permits their propagation to the next
block to be killed outside a larger ball of radius \(CL\) about
\(z_\alpha\), at a rapidly small relative loss of total variation.
Apply this restriction separately to each strand.  Apply the same
restriction about each injection site to the cutoff sources.  The
elapsed propagation time is at most \(4L\), also for the final block
continued into the copied period.

These balls lie in the unpinned interior.  Up to the first exit from
its chosen ball, the propagation conserves counting mass: the local
heat evolution, hole transports, and seam averages all do so, and
there is no pin killing there.  The localized propagation loses mass
only at that exit.  Its signed net mass therefore differs from the
signed sum of the injected sources by at most their total variation
times the rapidly small exit probability.  All source variations and the number of
injection sites are polynomially bounded in \(\beta\) on
\(\mathcal G\).  Their aggregate loss is consequently rapidly small,
uniformly in the volume and in the averaging slice.

Fix a fully regular averaging slice.  For each \(\alpha\), choose
one available representative slot \(a_\alpha\) within a bounded
distance of \(z_\alpha\).  Proposition~\ref{prop:geometryinputs}
provides this slot and paths from every slot carrying the localized
mass to \(a_\alpha\), of length and displacement at most \(CL\).
Erase loops from these paths.  Routing a signed mass along them
produces a flow whose divergence is that mass minus its net mass at
\(a_\alpha\).  Its norm is at most the total variation times the
largest path norm, hence its squared norm is bounded by
\[
 CL\bigl(LW(z_\alpha)\bigr)^2=CL^3W(z_\alpha)^2.
\]
At a given edge, only strands starting within \(CL\) can contribute.
There are at most \(CL^3\) such starting slots, by deterministic
counting and fixed slot multiplicity.  Cauchy--Schwarz on each edge
therefore bounds the squared norm of the sum of these local flows by
\begin{equation}\label{eq:blocklocalflow}
 CL^6Q_j.
\end{equation}

Let \(q_\alpha\) be the remaining net mass at \(a_\alpha\).
The cancellation in Lemma~\ref{lem:cancellation}, and the preceding
conservation observation, give
\[
 |q_\alpha|\le CL\theta_\alpha W(z_\alpha).
\]
Indeed the error from localization is rapidly small relative to
\(LW(z_\alpha)\), so for an isolated strand it is absorbed into
\(CL^{3/4}W(z_\alpha)\).  Send \(q_\alpha\) to infinity along
the unit flow from \(a_\alpha\).  Bounded displacement of these
representatives preserves the inverse-distance overlap estimate in
Lemma~\ref{lem:flowoverlap}, up to a constant.  The squared norm of
the resulting sum of flows is thus at most
\begin{equation}\label{eq:blocknetflow}
 C\sum_{\alpha,\alpha'}
 \frac{|q_\alpha q_{\alpha'}|}
 {1+|z_\alpha-z_{\alpha'}|_\infty}
 \le CL^2C_j.
\end{equation}
This is the only part of the obstacle flow that carries mass over arbitrarily
large distances, and it is the part on which isolated-strand
cancellation saves a factor \(L^{-1/2}\) after squaring.

For the cutoff sources, each injection site contributes total variation
at most \(CLb_*\) in the block.  Send their localized descendants
directly to infinity using Lemma~\ref{lem:flowoverlap}.  If two
injection sites are within \(CL\), use the constant overlap bound.
There are at most \(CR^3L^3\) such ordered pairs.  For sites farther
apart, the descendants are still separated by a constant multiple of
their starting distance, after increasing the threshold \(CL\).
The inverse-distance sum over a ball of radius \(CR\) is \(CR^5\).
Their total squared flow norm is at most
\begin{equation}\label{eq:blockcutoffflow}
 CL^2b_*^2(R^3L^3+R^5)\le CL^2b_*^2R^5.
\end{equation}
Adding the three flows and using the squared triangle inequality proves
\eqref{eq:blockflowbound}.  Every source of these flows is within
\(B_{6R}(x)\), since \(L=o(R)\).

We next average the bound over the exposure.  Radial shells give
\begin{equation}\label{eq:weightenergysums}
 \sum_{|z-x|_\infty\le4R}W(z)^2
 \le C\beta^{-2}D_*^{-1},
 \qquad
 \sum_{|z-x|_\infty,|z'-x|_\infty\le4R}
 \frac{W(z)W(z')}{1+|z-z'|_\infty}
 \le C\beta^{-2}R.
\end{equation}
For the second estimate one may enlarge \(W(z)\) to
\(\beta^{-1}(1+|z-x|_\infty)^{-2}\).
Two comparable dyadic shells of radius \(b\) contribute at most
\(C\beta^{-2}b\): for each site in one shell, the inverse-distance
sum over the other is at most \(Cb^2\).
Shells of radii \(c\le b/2\) contribute at most
\(C\beta^{-2}c\).  Summing first over \(c\) and then over the
dyadic values of \(b\le CR\) proves the estimate without a
logarithmic factor.

Proposition~\ref{prop:snapshot} bounds the expectation of \(Q_j\) by
\(C\lambda\sum W^2\).  To estimate \(C_j\), first take pairs
whose starting sites are at distance at most \(C_1L\), where
\(C_1\) is a sufficiently large fixed constant compared with the
isolation radius.  Charge only one obstacle, with probability
\(C\lambda\), and allow \(CL^3\) possible second slots.
The envelope \(W\) is comparable at these distances because
\(L=o(D_*)\).  This portion of \(L^2C_j\) is bounded in
expectation by
\[
 CL^5\lambda\sum W^2,
\]
which is absorbed into the \(L^6\) term in
\eqref{eq:blockflowexpectation}.

For a more distant pair, the two starting slots are distinct.  If
both are isolated, their occupation probability is at most
\(C\lambda^2\), and \(\theta_\alpha\theta_{\alpha'}=L^{-1/2}\).
If at least one is not isolated, a third distinct obstacle slot
must lie within distance \(c_2L\) of that member.  The separation
ensures that this third slot is different from the other member of
the pair.  A union bound over its \(CL^3\) possible positions and
the three-point snapshot bound give probability at most
\(CL^3\lambda^3\).  Applying \eqref{eq:weightenergysums} proves
\eqref{eq:blockflowexpectation}.  Only one-time snapshot probabilities
have been used; no independence between strands or between blocks
is needed.

Finally, at the terminal cut, the localized mass from one strand
has \(\ell^2\) norm at most its variation \(CLW(z_\alpha)\).
The mass from one cutoff injection site has norm at most \(CLb_*\).
Their supports have overlap at most \(CL^3\).  Cauchy--Schwarz
at each slot proves \eqref{eq:lastblocknorm}.
The displayed bounds use measurable finite snapshot sums; choosing
the flows measurably is unnecessary.
\end{proof}

\subsection{Averaging over the next block}

The flow bound controls a block at any fully regular slice of its
successor.  We now use these choices of slice to share a single
Dirichlet energy budget among the blocks.  The final block requires
a separate step, because its successor lies beyond the terminal cut.

Let \(m_j\) be the terminal mass produced by the first collection of
sources in \(I_j\).  For \(j<n\), at a slice \(t\in I_{j+1}\)
the block pairing is the propagated source mass paired with
\(Y(t)=T_{\beta,t}^*y\).  Lemma~\ref{lem:blockflow} and
\eqref{eq:flowpairing-energy} therefore apply directly.
For the final block, put \(m=m_n\) and use the exact identity
\begin{equation}\label{eq:lastblockcopy}
 \ip{m}{y}
 =\ip{m}{y-U_d^*y}+\ip{m}{U_d^*y}.
\end{equation}
The second term equals \(\ip{U_dm}{y}\).  It is represented by
continuing the source mass through a copy of the period, with
terminal test vector \(y\) at time \(2\beta\).  At its averaging
slice \(t\in\beta+I_1\), the test vector is consequently
\(T_{2\beta,t}^*y\).  The localized norm bound and
\eqref{eq:energy-suffix} control the first term:
\begin{equation}\label{eq:lastblockcorrection}
 \abs{\ip{m}{y-U_d^*y}}
 \le \sqrt{2K_ne_y}+C_0\rho(\beta).
\end{equation}
The localization error here is estimated in \(\ell^1\), using
\(\norm{y-U_d^*y}_\infty\le2C_0\).
Figure~\ref{fig:period-averaging} shows the ordinary averaging windows
and the window used for this continued final block.
\begin{figure}[tbp]
\centering
\begin{tikzpicture}[x=0.9cm,y=0.9cm,>=Stealth,
 every node/.style={font=\small}]
 \definecolor{flowblue}{RGB}{25,108,105}
 \foreach \a/\b/\lab in {0/1.3/I_1,2.5/3.8/I_j,3.8/5.1/I_{j+1},6.4/7.7/I_n}
  {\draw[fill=gray!7] (\a,0) rectangle (\b,0.75);
   \node at ({(\a+\b)/2},0.375) {$\lab$};}
 \node at (1.9,0.375) {$\cdots$};
 \node at (5.75,0.375) {$\cdots$};
 \draw[fill=flowblue!10] (7.7,0) rectangle (9.5,0.75);
 \node at (8.6,0.375) {$I_1+\beta$};
 \node at (10.05,0.375) {$\cdots$};
 \draw[->] (-0.25,-0.25)--(10.65,-0.25) node[right] {$t$};
 \draw (0,-0.2)--(0,-0.3) node[below] {$0$};
 \draw (7.7,-0.2)--(7.7,-0.3) node[below] {$\beta$};
 \draw[dashed] (7.7,-0.05)--(7.7,2.1);
 \draw[->,thick,flowblue] (3.15,0.9) to[out=70,in=110] (4.45,0.9);
 \node[align=center,text=flowblue] at (3.8,1.9) {propagate, then average\\on regular slices};
 \draw[->,thick,flowblue] (7.05,0.9) to[out=70,in=110] (8.6,0.9);
 \node[align=center,text=flowblue] at (8.5,2.55) {continue through\\the next period};
 \node[align=center] at (5.1,-1.3)
  {$\langle m,y\rangle=\langle m,y-U_d^*y\rangle+\langle U_dm,y\rangle$};
\end{tikzpicture}
\caption{Averaging the source mass from one block in the next block.
The last block uses the first block of a copied period. The displayed identity
separates the pairing of its terminal mass $m$ into a term controlled by contraction energy
and a term whose propagation can be split at the copied averaging slices.
The averaging windows are disjoint modulo one period, so their Dirichlet
energies are summed with a constant total cost. Block widths are schematic.}
\label{fig:period-averaging}
\end{figure}


For each \(j\), average over the fully regular subset of its
successor block.  Proposition~\ref{prop:good} gives this subset
measure at least \(cL\).  Applying Cauchy--Schwarz in that time
variable bounds its averaged localized pairing by
\begin{equation}\label{eq:blockaverage}
 C\sqrt{F_j/L}
 \left(\int_{I_{j+1}}\D_t(Y_j(t))\,dt\right)^{1/2},
\end{equation}
where \(Y_j(t)=T_{\beta,t}^*y\) for \(j<n\), whereas
\(Y_n(t)=T_{2\beta,t}^*y\) and \(I_{n+1}=\beta+I_1\).
The intervals \(I_2,\ldots,I_n\) are disjoint in the original
period, and \(\beta+I_1\) lies in the copied period.  Each full
period has the bound \eqref{eq:energy-suffix}, so
\[
 \sum_{j=1}^n\int_{I_{j+1}}\D_t(Y_j(t))\,dt\le2e_y.
\]
Cauchy--Schwarz across the block indices in
\eqref{eq:blockaverage}, followed by \eqref{eq:lastblockcorrection},
gives
\begin{equation}\label{eq:smoothdualcoefficient}
 \left|\sum_{j=1}^n\ip{m_j}{y}\right|
 \le C\left(\sum_{j=1}^nF_j/L+K_n\right)^{1/2}\sqrt{e_y}
       +C_0\rho(\beta).
\end{equation}
The sum of localization errors is still rapidly small because
\(n\le C\beta/L\).

The expectation of the first term in parentheses is bounded by
\(C\beta/L^2\) times the right side of
\eqref{eq:blockflowexpectation}.  Substituting
\(\lambda=\beta^{-3/2}\), \(D_*=\sqrt\beta\), and
\(b_*=\beta^{-1}R^{-3}\), this is
\begin{equation}\label{eq:smoothdualexpectation}
 C\left[
 L^4\beta^{-3}
 +L^{-1/2}\beta^{-4}R
 +L^3\beta^{-11/2}R
 +\frac1{\beta R}\right].
\end{equation}
Likewise, \eqref{eq:lastblocknorm} and the one-point snapshot bound
give
\begin{equation}\label{eq:copycorrectionexpectation}
 \EE_P[\ind_{\mathcal G}K_n]
 \le C\left(L^5\beta^{-4}+L^5\beta^{-2}R^{-3}\right).
\end{equation}
Thus the block argument has converted cancellation of source mass
into a dual energy bound.  It remains to handle the early
point-profile sources, for which cancellation at the scale \(L\)
was not used.

\subsection{Early point-profile sources}

\begin{lemma}[Diffusion of the early sources]\label{lem:earlysourceenergy}
Let \(\mu_d^{\mathrm{early}}\) be the terminal mass of the
point-profile obstacle sources at elapsed times in \([0,\beta/2)\),
and let \(M\) be their total variation before propagation.  Then
\begin{equation}\label{eq:earlymassmoment}
 \EE_P[\ind_{\mathcal G}M^2]\le H_\beta\beta^{-3/2}.
\end{equation}
For every test vector of Proposition~\ref{prop:dualenergy}, on
\(\mathcal G\),
\begin{equation}\label{eq:earlydualbound}
 \abs{\ip{\mu_d^{\mathrm{early}}}{y}}
 \le CM\beta^{-3/4}\sqrt{e_y}+C_0\rho(\beta).
\end{equation}
\end{lemma}

\begin{proof}
For a bin beginning at elapsed time \(s\), put
\(d_s=\sqrt{1+s}\).  Lemma~\ref{lem:sourcebins} bounds its source
variation on \(\mathcal G\) by
\(H_\beta\sum_\alpha W_{d_s}(z_\alpha)\), where the sum is over
the obstacle slots above \(B_{4R}(x)\) at that bin's starting slice.
The single-box
case of Lemma~\ref{lem:boxmoment} gives
\begin{align*}
 \left\|\sum_\alpha W_{d_s}(z_\alpha)\right\|_{L^2(\PP_P)}
 &\le C\left(\sqrt\lambda\,\norm{W_{d_s}}_2
                +\lambda\norm{W_{d_s}}_1\right)\\
 &\le C\left(\beta^{-3/4}d_s^{-5/2}
                 +\beta^{-3/2}d_s^{-1}\right)=J(d_s).
\end{align*}
Minkowski's inequality in the probability space now gives
\[
 \|\ind_{\mathcal G}M\|_{L^2(\PP_P)}
 \le H_\beta\sum_{\mathrm{bins}}J(d_s)
 \le H_\beta\beta^{-3/4}.
\]
For the last inequality, bins have lengths between \(1\) and \(3\),
so the sum of \(d_s^{-5/2}=(1+s)^{-5/4}\) is bounded and the sum of
\(d_s^{-1}=(1+s)^{-1/2}\) is at most \(C\sqrt\beta\).
Absorbing the square of a logarithmic factor into \(H_\beta\)
proves \eqref{eq:earlymassmoment}.  This use of Minkowski requires
no independence of the bin snapshots.

Propagate all these sources to a slice in a block contained in the
last quarter of the traversal.  Every propagation has elapsed time
at least \(\beta/4\).  Localize to \(B_{6R}(x)\); by
Lemma~\ref{lem:localization}, the discarded total variation is
rapidly small on \(\mathcal G\).  Lemma~\ref{lem:heatbound}
and substochasticity show that a localized unit point mass has total
mass at most one and maximum entry at most \(C\beta^{-3/2}\).
The additional spatial restriction only decreases the kernel.
Consequently the localized signed source mass \(\nu_t\) is
dominated in absolute value by a nonnegative mass \(\sigma_t\)
such that
\[
 \norm{\sigma_t}_1\le M,
 \qquad
 \norm{\sigma_t}_\infty\le CM\beta^{-3/2}.
\]

At a fully regular such slice, use the unit flows of
Lemma~\ref{lem:flowoverlap}.  For every slot \(a\), splitting
distances at \(D_*\) gives
\begin{align*}
 \sum_{a'}\frac{\sigma_t(a')}{1+\dist(a,a')}
 &\le CM\beta^{-3/2}D_*^2+CM D_*^{-1}\\
 &\le CM\beta^{-1/2}.
\end{align*}
Superposing the unit flows with coefficients \(\nu_t(a)\) therefore
gives a flow with divergence \(\nu_t\) and squared norm at most
\[
 C\sum_{a,a'}\frac{\sigma_t(a)\sigma_t(a')}
                         {1+\dist(a,a')}
 \le CM^2\beta^{-1/2}.
\]
There are blocks of total length at least \(c\beta\) contained
in the last quarter, and at least half of each is fully regular.
Average \eqref{eq:flowpairing-energy} over this regular set.
Cauchy--Schwarz in time and \eqref{eq:energy-suffix} give
\eqref{eq:earlydualbound}.
\end{proof}

\begin{proof}[Proof of Proposition~\ref{prop:dualenergy}]
The two collections exhaust the sources in the chosen direction.
Combine \eqref{eq:smoothdualcoefficient} and
\eqref{eq:earlydualbound}, and choose a measurable coefficient with
\[
 X_d^2\le C\left(\sum_{j=1}^nF_j/L+K_n+M^2\beta^{-3/2}\right).
\]
All expressions on the right are measurable functions of the
exposure; none depends on the test vector.  Equations
\eqref{eq:smoothdualexpectation},
\eqref{eq:copycorrectionexpectation}, and
\eqref{eq:earlymassmoment} bound its expectation on
\(\mathcal G\).  With \(L=\beta^\delta\) and
\(R=\beta^{3/2+\delta/4}\), their sum is at most
\begin{equation}\label{eq:energyfinalpowers}
 C\left[
 \beta^{-3+4\delta}
 +\beta^{-5/2-\delta/4}
 +\beta^{-4+13\delta/4}
 +\beta^{-4+5\delta}
 +\beta^{-13/2+17\delta/4}
 +H_\beta\beta^{-3}\right]
 =o(\beta^{-5/2}),
\end{equation}
since \(\delta=1/200\).
The term \(\beta^{-5/2-\delta/4}\) covers both the isolated
net-mass flows and the cutoff flows.  All localization errors were
estimated in \(\ell^1\) against test vectors of supremum norm
at most a constant times \(C_0\), so they give the uniform
rapidly small remainder in \eqref{eq:dualenergy}.
The argument used elapsed time and the ordered kernels throughout;
it applies to \(U\) and to \(U^*\) with the same bounds.
Adding the two expected squares proves
\eqref{eq:dualenergyexpectation}.
\end{proof}
